\sheet{12}{Knock Detection}{Chapter~12 (capstone: angle-synchronous knock detection)}

\begin{goals}
\begin{itemize}
  \item Assemble the methods of the previous sheets into a complete,
    angle-synchronous knock detector and justify every block quantitatively.
  \item Evaluate a detector statistically on thousands of labelled combustions:
    ROC curve, area under the curve, detection rate at a fixed false-alarm rate,
    choice of the threshold.
  \item Compare energy-, FFT- and MEM-based knock features and quantify the
    contribution of band-pass, angle window and background normalisation.
  \item Run the detector in real time on the ESP32 with crank synchronisation,
    per-cylinder knock indication and ignition retard; optionally port it to the
    8-bit Uno in fixed point.
\end{itemize}
\end{goals}

\section*{The signal chain}

\ecuquote{} The ECU never sees the knock -- only a vibration signal that also
contains noise, the combustion vibration and valve impacts. The standard knock
detector of a production ECU is the chain below; every block is one of the
previous sheets.

\begin{center}
\begin{tikzpicture}[>=Stealth,font=\footnotesize,node distance=3mm,
  blk/.style={draw,rounded corners=1pt,align=center,minimum height=9mm,inner sep=2pt,fill=dspblue!6}]
  \node[blk] (adc) {ADC\\$f_s$=25\,kHz};
  \node[blk,right=of adc] (bp) {band-pass\\6--7\,kHz};
  \node[blk,right=of bp] (sq) {$(\cdot)^2$};
  \node[blk,right=of sq] (win) {window\\$10^\circ$--$70^\circ$};
  \node[blk,right=of win] (int) {$\sum$};
  \node[blk,right=of int] (nrm) {$\div B_c$};
  \node[blk,right=of nrm] (thr) {$>K$?};
  \node[blk,below=6mm of win] (sync) {crank sync\\(60-2, CAM)};
  \node[blk,right=of thr] (act) {retard\\LED};
  \node[blk,below=7mm of nrm,xshift=8mm] (bg) {background\\$B_c$ (EWMA)};
  \draw[->] (adc) -- (bp); \draw[->] (bp) -- (sq); \draw[->] (sq) -- (win);
  \draw[->] (win) -- (int); \draw[->] (int) -- (nrm); \draw[->] (nrm) -- (thr);
  \draw[->] (thr) -- (act); \draw[->] (sync) -- (win);
  \draw[->] (int.south) |- (bg.west); \draw[->] (bg) -- (nrm);
\end{tikzpicture}
\end{center}

Conventions: $f_s=\SI{25}{kHz}$ (Sheet~10), band-pass = the 10th-order Butterworth
biquad cascade of \texttt{knock\_filter.h} (Sheet~10; group delay 22 samples at
\SI{6.5}{kHz}) or the single biquad \texttt{kf\_bp2} ($B=\SI{400}{Hz}$, 20 samples),
window $10^\circ$--$70^\circ$ after the TDC of the firing cylinder, feature
$E_c=\sum_{\text{window}}y^2[n]$ for cylinder $c$, normalised
$R_c=E_c/B_c$, knock if $R_c>K$. Ground truth: \texttt{e.knock\_flag[c]} in the
C model, the KNK pin of the engine-sim chip in Wokwi. Code:
\codefile{jslinux/sheet12/} (\texttt{knockeval.c} plus \texttt{knock\_filter.h},
\texttt{ar.h}, \texttt{fft.h} -- use your own versions from Sheets 8, 10, 11 if you
like) and \codefile{wokwi/sheet12\_knock\_esp32/}, \codefile{wokwi/sheet12\_knock\_uno/}.

%-------------------------------------------------------------------------
\begin{exercise}{Detector design}{\Paper}
\begin{tasks}
  \item Justify every block of the chain with the respective chapter: why
    sampling at \SI{25}{kHz} (and what about anti-aliasing), why an angle
    window and why exactly $10^\circ$--$70^\circ$ ATDC, why a band-pass,
    why squaring (instead of a linear operation), why summation, why a
    normalisation and why a threshold. Which blocks are LTI, which are not?
  \item Compute the window length in samples at 1500, 3000 and \SI{6000}{rpm}.
    By how many degrees must the window be shifted to compensate the group delay
    of the band-pass (22 samples)? Why is a fixed sample shift correct although the
    window is defined in degrees?
  \item Energy budget. Show that a knock burst $A\e^{-t/\tau}\sin(2\pi f_1t)$ that lies
    completely in the window and is passed with gain 1 contributes
    $E_k\approx f_sA^2\tau/4$ (in V$^2$, sum over samples). White noise with standard
    deviation $\sigma$ contributes on average $E_n=N_w\sigma^2\cdot2B_n/f_s$, where
    $N_w$ is the window length and $B_n$ the noise bandwidth of the band-pass
    ($B_n=\int_0^{f_s/2}|H(f)|^2\dd f=\SI{1161}{Hz}$ for the cascade). Evaluate $E_k$ for
    $A=\SI{0.2}{V}$ and $E_n$ for $\sigma=0.05$ and \SI{0.15}{V} at the three speeds.
    What do you expect for the detection performance?
  \item Without knock, $E_c$ is approximately $E_n\cdot\chi^2_\nu/\nu$ distributed
    with $\nu\approx2B_nT_w$ degrees of freedom ($T_w$ = window duration). Compute
    $\nu$ for the three speeds and, using the 99\,\% quantiles
    $\chi^2_{\nu,0.99}/\nu\approx 2.0$ ($\nu=15.5$), $2.6$ ($\nu=7.7$),
    $3.4$ ($\nu=3.9$), the threshold for 1\,\% false alarms in units of the mean noise
    energy. Why does a fixed threshold $K$ give an rpm-dependent false-alarm rate?
  \item Background normalisation: propose an update rule for $B_c$ that follows slow
    changes (sensor ageing, rpm, load) but is not pulled up by knocking combustions.
    Why one background per cylinder?
  \item Control action: on a detected knock the ignition of that cylinder is retarded by
    $\Delta_r$, without knock it is advanced again by $\Delta_a\ll\Delta_r$ per
    combustion. Explain the asymmetry and what a false alarm costs.
\end{tasks}
\end{exercise}

\begin{solution}
a) \emph{Sampling} (Ch.~4/5): the knock modes (6.5/10.5\,kHz) need $f_s>\SI{21}{kHz}$;
\SI{25}{kHz} fits the \SI{50}{kHz} update of the chip (Sheet~10). A real ECU needs an
analog anti-aliasing low-pass below \SI{12.5}{kHz}, otherwise broadband vibration above
$f_s/2$ folds into the knock band.
\emph{Crank sync and angle window} (Ch.~1, 4): knock can only occur after the spark,
during the combustion of one cylinder, typically 5--30$^\circ$ ATDC, and decays within
a few ms; the window excludes the noise of the rest of the cycle and the valve impacts
at $90^\circ$ ATDC (and, in a real engine, injector and piston-slap noise). Since the
events are fixed in \emph{angle}, the window must be angle based (time-based windows
would move with rpm). $10^\circ$ start: the combustion vibration and spark noise before;
$70^\circ$ end: before the valve impact.
\emph{Band-pass} (Ch.~10): selects the knock resonance and rejects the low-frequency
combustion vibration (\SI{0.2}{V}, much larger than light knock) and most of the
broadband noise (gain in SNR $\approx (f_s/2)/B_n\approx 11$, i.e.\ \SI{10}{dB}).
\emph{Squaring} (Ch.~2, nonlinear memoryless system): the knock oscillation has
random phase, its linear average is zero; the energy (or the rectified amplitude) is
phase independent. \emph{Summation} (Ch.~3, LTI, moving sum): averages the fluctuations
of $y^2$ over the window (matched to the burst duration).
\emph{Normalisation}: the absolute energy depends on rpm (window length), sensor
sensitivity, mounting position, engine load and ageing; the ratio to the background is
dimensionless and comparable. \emph{Threshold}: a binary decision (Ch.~1: from signal to
information) with a trade-off between misses and false alarms.
LTI: band-pass, summation (with a fixed window). Not LTI: windowing (time-variant,
angle-synchronous), squaring (nonlinear), normalisation and threshold (nonlinear,
adaptive).

b) $f_s=\SI{25}{kHz}$, $60^\circ$ take $60/(6\cdot\text{rpm})$\,s:
\SI{6.67}{ms} = 167 samples at 1500, \SI{3.33}{ms} = 83 at 3000, \SI{1.67}{ms} = 42 at
\SI{6000}{rpm}. The delay of $22T_s=\SI{0.88}{ms}$ corresponds to
$7.9^\circ$, $15.8^\circ$, $31.7^\circ$. The filter delays the signal by a fixed
\emph{time}: the output sample $y[n]$ belongs to the input at $n-22$; hence the window
decision for $y[n]$ must use the crank angle of sample $n-22$ (in code: a ring buffer of
the last angles). Expressed in degrees the shift grows with rpm.

c) $\sum_n y^2[n]\approx f_s\int_0^\infty A^2\e^{-2t/\tau}\sin^2(2\pi f_1t)\dd t
=f_sA^2\cdot\tfrac{\tau}{2}\cdot\tfrac12=f_sA^2\tau/4=5A^2$ (V$^2$) $\Rightarrow$
$E_k=\SI{0.20}{V^2}$ for $A=\SI{0.2}{V}$ (the knock amplitude of the model varies
between 0.5 and 1.5 times this value, i.e.\ $E_k=0.05\dots0.45$, and part of a late
burst falls outside the window). Noise: $E_n=N_w\sigma^2\cdot2\cdot1161/25000$:
\begin{center}
\small
\begin{tabular}{@{}rrrr@{}}
\toprule
rpm & $N_w$ & $E_n$, $\sigma=\SI{0.05}{V}$ & $E_n$, $\sigma=\SI{0.15}{V}$\\
\midrule
1500 & 167 & 0.039 & 0.348\\
3000 & 83 & 0.019 & 0.174\\
6000 & 42 & 0.010 & 0.087\\
\bottomrule
\end{tabular}
\end{center}
At $\sigma=\SI{0.05}{V}$ the knock energy is 5--20 times the mean noise energy: good
detection, better at high rpm (shorter window, less noise). At $\sigma=\SI{0.15}{V}$
the ratio is only 0.6--2.3: weak knocks are buried in the noise, the detection rate at
low false-alarm rates will be poor, especially at low rpm.

d) $\nu=2B_nT_w=15.5$, 7.7, 3.9 at 1500, 3000, \SI{6000}{rpm}; threshold for 1\,\% false
alarms $\approx2.0$, 2.6, $3.4\times$ the mean noise energy. With few degrees of freedom
(short window) the noise energy fluctuates more, the $\chi^2$ distribution has a heavier
tail, and a higher threshold is needed for the same false-alarm rate. A fixed $K$ is
therefore conservative at low rpm and produces more false alarms at high rpm; production
ECUs use rpm-dependent thresholds (maps). Note that normalisation removes the
\emph{mean} dependence on rpm but not the \emph{shape} of the distribution.

e) Exponentially weighted moving average with clipping:
$B_c\leftarrow B_c+\lambda\bigl(\min(E_c,\,3B_c)-B_c\bigr)$, $\lambda\approx0.05$ (time
constant $\approx20$ combustions of that cylinder, i.e.\ about 1.6\,s at
\SI{3000}{rpm}), larger $\lambda$ during the start (warm-up). The clipping limits the
influence of a knocking combustion (alternative: update only with combustions classified
as non-knocking -- but then a wrong decision can lock the background). One background per
cylinder because the transfer path from each cylinder to the sensor differs (distance,
block structure): the same knock intensity gives different sensor amplitudes.

f) Knock can damage the engine within few cycles -- fast reaction ($\Delta_r\approx
1.5^\circ$ per event). Retard costs efficiency and torque, so the ignition is advanced
again slowly ($\Delta_a\approx0.05^\circ$ per combustion) towards the knock limit; the
controller then oscillates closely below the limit. A false alarm costs efficiency
($\Delta_r$ that is recovered only after $\Delta_r/\Delta_a=30$ combustions of that
cylinder); a miss risks damage. Hence the threshold is chosen for a low but non-zero
false-alarm rate (e.g.\ 1\,\%).
\end{solution}

%-------------------------------------------------------------------------
\begin{exercise}{Offline evaluation with ground truth}{\JSLinux}
\sloppy
\texttt{knockeval.c} runs the virtual engine at 1500, 3000 and \SI{6000}{rpm} with sensor
noise $\sigma=0.05$ and \SI{0.15}{V} (6 scenarios, 1000 evaluated combustions each, knock
probability 0.2, knock amplitude \SI{0.2}{V}$\times(0.5\dots1.5)$, light knock) and computes
for every combustion eight features, each normalised by its own per-cylinder background:
\begin{center}
\small
\begin{tabularx}{\textwidth}{@{}lX@{}}
\toprule
WB-seg / WB-win & wide-band energy $\sum x^2$ over the whole $180^\circ$ segment / in the window\\
BP-seg / BP-win & band-pass (cascade) energy over the segment / in the delay-compensated window\\
BP2-win & as BP-win with the single biquad ($B=\SI{400}{Hz}$)\\
FFT / FFT-rect & power 6--7\,kHz of the zero-padded FFT of the raw window samples, Hann / no window\\
MEM & power 6--7\,kHz of the Burg AR(8) spectrum of the raw window samples\\
\bottomrule
\end{tabularx}
\end{center}
\begin{lstlisting}[style=shell]
gcc -O2 -o knockeval knockeval.c -lm && ./knockeval [knock_amp] [ncomb]
\end{lstlisting}
\begin{tasks}
  \item Complete the window logic: integrate $x^2$ in the window $10^\circ$--$70^\circ$
    ATDC and $y^2$ (cascade, single biquad) while the \emph{delayed} angle is in the
    window.
  \item Complete the background update (clipped EWMA, faster during the first 16
    combustions of each cylinder).
  \item Complete the ROC computation in \texttt{roc()}: the combustions are sorted by
    decreasing feature value; every threshold between two values gives a point
    (FPR, TPR). Compute the AUC (trapezoidal rule) and the detection rate at
    $\le1\,\%$ false alarms.
  \item Complete the FFT and MEM features (\texttt{fft.h}, \texttt{ar.h}).
  \item Run the program and interpret the AUC and detection tables: what is the
    contribution of the band-pass, of the angle window, of the delay compensation, of
    the normalisation (compare the pooled ROC with and without)? Plot \texttt{roc.csv}.
  \item Choose a threshold $K$ for 1\,\% false alarms over all scenarios and compare the
    per-scenario thresholds with your prediction from 12.1\,d).
  \item Compare the energy, FFT and MEM features. Why is the Hann-windowed FFT so much
    worse at \SI{1500}{rpm} than the rectangular one? Which feature would you implement
    on the microcontroller, and why?
\end{tasks}
\end{exercise}

\begin{solution}
a)--d) Window logic and background update
(\codefile{solutions/jslinux/sheet12/knockeval.c}):
\inputcode[firstline=163,lastline=189]{solutions/jslinux/sheet12/knockeval.c}
ROC (after sorting by decreasing value; ties are merged):
\inputcode[firstline=220,lastline=230]{solutions/jslinux/sheet12/knockeval.c}

e) Output (6000 combustions, 1210 with knock = 20.2\,\%):
\begin{center}
\small
\setlength{\tabcolsep}{4pt}
\begin{tabular}{@{}rr|rrrrrrrr@{}}
\toprule
\multicolumn{2}{c|}{AUC} & WB-seg & WB-win & BP-seg & BP-win & BP2-win & FFT & FFT-rect & MEM\\
\midrule
1500 & 0.05 & 0.799 & 0.788 & 0.957 & \textbf{0.981} & 0.978 & 0.804 & 0.973 & 0.974\\
1500 & 0.15 & 0.600 & 0.619 & 0.738 & \textbf{0.789} & 0.784 & 0.620 & 0.764 & 0.761\\
3000 & 0.05 & 0.811 & 0.819 & 0.985 & \textbf{0.998} & 0.996 & 0.955 & 0.997 & 0.997\\
3000 & 0.15 & 0.606 & 0.652 & 0.747 & \textbf{0.828} & 0.818 & 0.717 & 0.824 & 0.818\\
6000 & 0.05 & 0.912 & 0.922 & 0.988 & \textbf{0.999} & 0.998 & 0.997 & 0.999 & 0.999\\
6000 & 0.15 & 0.698 & 0.726 & 0.822 & \textbf{0.886} & 0.864 & 0.826 & 0.878 & 0.858\\
\multicolumn{2}{r|}{pooled} & 0.740 & 0.752 & 0.897 & \textbf{0.934} & 0.927 & 0.846 & 0.927 & 0.924\\
\midrule
\multicolumn{2}{c|}{det.\ @1\,\% FA} &&&&&&&&\\
1500 & 0.05 & 20.7 & 26.4 & 77.7 & 82.9 & 82.4 & 34.2 & 80.3 & 80.8\\
1500 & 0.15 & 3.5 & 2.0 & 7.0 & 20.4 & 14.4 & 5.5 & 15.4 & 18.9\\
3000 & 0.05 & 29.7 & 28.2 & 80.9 & 95.7 & 90.0 & 66.5 & 92.8 & 91.9\\
3000 & 0.15 & 2.1 & 3.7 & 16.8 & 30.9 & 24.6 & 14.1 & 27.7 & 24.6\\
6000 & 0.05 & 52.6 & 56.5 & 87.1 & 96.2 & 94.3 & 93.3 & 97.1 & 96.7\\
6000 & 0.15 & 4.3 & 8.7 & 29.0 & 37.2 & 34.8 & 24.6 & 31.9 & 24.6\\
\multicolumn{2}{r|}{pooled} & 11.8 & 9.3 & 44.5 & 59.1 & 55.9 & 40.3 & 57.1 & 56.4\\
\bottomrule
\end{tabular}
\end{center}
\emph{Band-pass:} the decisive block. Without it (WB) the combustion vibration and the
noise of all frequencies dominate; AUC 0.74--0.75 pooled, only 9--12\,\% detection at
1\,\% false alarms. With it (BP-seg) the AUC rises to 0.90.
\emph{Angle window:} BP-win vs.\ BP-seg: AUC 0.897 $\to$ 0.934, detection 44.5\,\%
$\to$ 59.1\,\%; integrating over $180^\circ$ collects three times more noise energy
(and the valve burst, which however is strongly attenuated by the cascade, \SI{-44}{dB}
at \SI{8}{kHz}). For the wide-band features the window helps little because the
low-frequency combustion vibration lies inside the window.
\emph{Delay compensation:} included in BP-win; without it, at \SI{6000}{rpm} a
$32^\circ$ portion of the burst would be lost (cf.\ Sheet~10, 10.3).
\emph{Normalisation:} pooled over all scenarios the raw BP-win energy reaches only
AUC 0.812 and 9.8\,\% detection at 1\,\% false alarms, the normalised ratio 0.934 and
59.1\,\%: the mean noise energy differs by a factor of 36 between the scenarios
(Table in 12.1\,c), so no single absolute threshold fits.
\emph{Noise:} at $\sigma=\SI{0.15}{V}$ even the best feature detects only 20--37\,\% of
the (light) knocks at 1\,\% false alarms -- consistent with the energy budget of 12.1\,c).
The ROC curves (\texttt{roc.csv}) of BP-win lie far above those of WB-win; at FPR = 5\,\%
the pooled TPR of BP-win is 74\,\% (WB-win: 28\,\%).

f) Pooled threshold for 1\,\% false alarms: $R>2.00$ (detection 59.1\,\%) -- this is the
value $K=2$ used in 12.3. Per scenario: $R>1.57$, 1.54, 2.02 ($\sigma=\SI{0.05}{V}$;
1500/3000/\SI{6000}{rpm}) and 1.79, 2.04, 2.49 ($\sigma=\SI{0.15}{V}$). The mean ratio of the
non-knocking combustions is below 1 (0.54--0.91), because the background contains
(clipped) knock energy; in units of the mean noise energy the thresholds are 2.3/2.0
(1500), 2.7/2.4 (3000) and 3.7/3.2 (\SI{6000}{rpm}) -- in good agreement with the
$\chi^2$ prediction 2.0, 2.6, 3.4 of 12.1\,d): the shorter the window, the higher the
threshold must be for the same false-alarm rate.

g) The band-pass energy (cascade) is the best and cheapest feature; the single biquad
(BP2) is almost as good (AUC 0.927 vs.\ 0.934): the matched 400-Hz band-pass rejects
noise as well as the steep cascade -- the Uno version is therefore not much worse.
FFT-rect and MEM are comparable to BP-win at $\sigma=\SI{0.05}{V}$ and slightly worse at
high noise; they cost much more (FFT of 256 points, or Burg + spectrum evaluation, per
combustion) and give no extra information here, since there is only one knock mode in
the band. The Hann-windowed FFT is poor at \SI{1500}{rpm} (AUC 0.80 vs.\ 0.97): the
window is \SI{6.7}{ms} long, but the burst starts at the beginning of the window
(onset $5^\circ$--$25^\circ$ ATDC) and has decayed after $\approx\SI{2}{ms}$ -- exactly
where the Hann window is close to zero. At \SI{6000}{rpm} (window \SI{1.7}{ms}) the burst
fills the window and the difference vanishes. Spectral features become interesting if
several knock modes must be weighted, if the knock frequency varies, or to distinguish
knock from disturbances at nearby frequencies with very short windows (MEM, Sheet~11).
Recommendation for the microcontroller: band-pass energy in the delay-compensated angle
window, normalised per cylinder.

\emph{Limitation of the model:} the valve impacts and the combustion vibration of the
virtual engine are deterministic (same amplitude every cycle); a constant disturbance
only shifts the feature and is largely removed by the normalisation. In a real engine
these disturbances vary from cycle to cycle and the angle window becomes even more
important.
\end{solution}

%-------------------------------------------------------------------------
\begin{exercise}{Real-time knock detector}{\Wokwi}
\sloppy
Open \codefile{wokwi/sheet12\_knock\_esp32/} (ESP32 + engine-sim chip, 4 knock LEDs,
retard LED, logic analyzer). Pins: KNOCK $\to$ GPIO34, CRANK $\to$ 18, CAM $\to$ 19,
TDC $\to$ 21 (only to check the synchronisation), KNK $\to$ 22 (ground truth, only for the
evaluation); LEDs cylinder 1--4 on 25, 26, 27, 14; retard PWM on 13. The sketch samples at
\SI{25}{kHz} (deadline loop), filters with \texttt{kf\_sos} and prints one line
\texttt{cyl,rpm,R,knock,truth,retard} per combustion and every 200 combustions the
confusion counts.
\begin{tasks}
  \item \emph{Crank sync.} Complete the 60-2 decoder in the CRANK interrupt (gap
    detection: a tooth period more than twice the previous one; revolution from the CAM
    level at tooth 10) and the angle interpolation \texttt{crank\_angle()}. The sketch
    compares the decoded angle at every TDC pulse with 120/300/480/660$^\circ$ and prints
    the mean error. What limits the accuracy?
  \item \emph{Window.} Complete the integration of $y^2$ while the delayed angle (ring
    buffer \texttt{th\_hist}, \texttt{KF\_DELAY} = 22 samples) is in the window.
  \item \emph{Decision and action.} Complete \texttt{finish\_combustion()}:
    $R=E/B_c$, knock if $R>K$ ($K=2$ from 12.2), background update, ignition retard
    ($+1.5^\circ$ per knock, $-0.05^\circ$ per knock-free combustion, max.\ $9^\circ$).
  \item Run the detector at 1500, 3000, \SI{6000}{rpm} with knock amplitude 1.0 and
    \SI{0.2}{V} and noise 50 and \SI{150}{mV} (knock probability 20\,\%). Report detection
    and false-alarm rates and compare them with 12.2. Watch the LEDs and the retard
    output.
  \item Estimate the processing time per sample on the ESP32 and the UART load at
    \SI{6000}{rpm}. Why must the sketch not print the samples?
  \item The engine-sim chip has no spark input: the retard does not influence the knock
    probability. Describe how a closed knock control loop would behave and how you would
    test it.
\end{tasks}
\end{exercise}

\begin{solution}
a) Decoder (\codefile{solutions/wokwi/sheet12\_knock\_esp32/sketch.ino}):
\inputcode[firstline=53,lastline=91]{solutions/wokwi/sheet12_knock_esp32/sketch.ino}
After the first gap (at most one revolution) the tooth counter is valid; the
revolution flag is toggled at every gap and confirmed by the CAM level at tooth 10
(CAM HIGH for $0^\circ$--$180^\circ$). Mean TDC sync error (PC test of the sketch with the chip
model): $\le0.01^\circ$ at all speeds -- the TDC pulses coincide with the rising edges of
teeth 20 and 50, so this check verifies the tooth counting and the phase (a wrong
revolution flag would give $360^\circ$). Note that the TDC flag must be discarded while the
decoder is not yet synchronised (stale pulse). In Wokwi the interrupt latency (a few
\si{\micro s}, i.e.\ $\approx0.1^\circ$ at \SI{6000}{rpm}) adds to this. Between the teeth the
angle is interpolated with the last tooth period; the 1\,\% second-order speed ripple
changes the speed by only $\approx0.2\,\%$ within one tooth, so the interpolation error
stays below $0.01^\circ$; the largest error occurs over the gap ($18^\circ$ extrapolated).

b), c) Window and decision:
\inputcode[firstline=209,lastline=245]{solutions/wokwi/sheet12_knock_esp32/sketch.ino}
\inputcode[firstline=141,lastline=160]{solutions/wokwi/sheet12_knock_esp32/sketch.ino}

d) Results (PC test harness running the sketch against the engine-sim chip code;
counts after the warm-up; Wokwi results differ in detail):
\begin{center}
\small
\begin{tabular}{@{}rrrrrrrrr@{}}
\toprule
rpm & knock amp. & noise & combustions & TP & FP & FN & detection & false alarms\\
\midrule
3000 & \SI{1.0}{V} & \SI{50}{mV} & 600 & 115 & 1 & 0 & 100\,\% & 0.2\,\%\\
6000 & \SI{1.0}{V} & \SI{50}{mV} & 1400 & 279 & 4 & 0 & 100\,\% & 0.4\,\%\\
1500 & \SI{0.2}{V} & \SI{50}{mV} & 200 & 32 & 0 & 12 & 72.7\,\% & 0.0\,\%\\
3000 & \SI{0.2}{V} & \SI{50}{mV} & 600 & 104 & 1 & 11 & 90.4\,\% & 0.2\,\%\\
6000 & \SI{0.2}{V} & \SI{50}{mV} & 1400 & 267 & 6 & 12 & 95.7\,\% & 0.5\,\%\\
1500 & \SI{0.2}{V} & \SI{150}{mV} & 200 & 4 & 1 & 40 & 9.1\,\% & 0.6\,\%\\
3000 & \SI{0.2}{V} & \SI{150}{mV} & 600 & 36 & 10 & 79 & 31.3\,\% & 2.1\,\%\\
6000 & \SI{0.2}{V} & \SI{150}{mV} & 1400 & 131 & 30 & 148 & 47.0\,\% & 2.7\,\%\\
\bottomrule
\end{tabular}
\end{center}
The real-time detector reproduces the offline results of 12.2 (same chain, same $K$):
strong knock is always detected; light knock (\SI{0.2}{V}) at \SI{50}{mV} noise in
73--96\,\% of the cases with $<1\,\%$ false alarms, and at \SI{150}{mV} noise the
detection collapses while the false-alarm rate rises above the design value at high
rpm (fewer degrees of freedom, 12.1\,d). The LED of a cylinder flashes for \SI{100}{ms}
at each detection, the retard LED brightness follows the mean retard (sawtooth: fast
up, slow down).

e) Per sample: \texttt{analogRead}, 5 biquads (25 multiply-adds in float), angle
interpolation (one division), a few comparisons -- some hundred instructions, i.e.\
$\approx1$--\SI{2}{\micro s} at \SI{240}{MHz} plus the ADC conversion, well below
\SI{40}{\micro s} (Wokwi timing is only indicative). UART: at \SI{6000}{rpm} 200
combustions/s $\times\approx22$ characters = 4400 characters/s, 38\,\% of the
11\,520 characters/s at 115\,200 baud; the 4-kB TX buffer decouples the printing from
the sampling loop. Printing the raw samples (25\,000/s $\times\ge5$ characters) would
exceed the UART capacity tenfold.

f) In a closed loop the retard lowers the cylinder pressure and temperature and thus the
knock probability; the controller then settles into a limit cycle around the knock limit
(retard steps after each knock, slow advance in between). Testing requires a knock
model that depends on the spark advance, e.g.\ $p_\text{knock}=f(\text{spark\_adv}-\text{retard})$,
i.e.\ a SPARK input or a serial/I$^2$C command input of the chip. One would then record
retard and knock rate per cylinder over time and check stability (no persistent heavy
knock, small mean retard) for different noise levels and false-alarm rates.
\end{solution}

%-------------------------------------------------------------------------
\begin{exercise}{Port to the Arduino Uno}{\Challenge}
\sloppy
Port the detector to the Uno (\codefile{wokwi/sheet12\_knock\_uno/}, starting from the Q15
biquad of Sheet~10): Timer1 ISR at \SI{25}{kHz}, single biquad \texttt{kf\_bp2\_q14},
window from the TDC pulse (period in samples of the previous combustion, cylinder~1 =
TDC while CAM is HIGH), energy $\sum (y^2\gg8)$ in 32 bit, integer background
$B\mathrel{+}=(\min(E,3B)-B)\gg4$, decision $E>2B$ without division. Report detection and
false-alarm rates as in 12.3 and the limits of the port.
\end{exercise}

\begin{solution}
See \codefile{solutions/wokwi/sheet12\_knock\_uno/sketch.ino}. The ISR adds only the
window bookkeeping and the KNK check (\texttt{PIND}) to the ISR of Sheet~10 (estimated
$\approx500$ of 640 cycles per sample); the decision, background, LEDs and printing run
in \texttt{loop()} once per combustion. Results of the PC test (knock probability 20\,\%):
1.0\,V/50\,mV at \SI{3000}{rpm}: 107 TP, 5 FP, 0 FN, 488 TN;
0.2\,V/50\,mV: 87\,\% detection / 0\,\% false alarms at 1500 (34/39), 92.5\,\% / 1.2\,\% at
3000 (99/107), 91.9\,\% / 0.1\,\% at \SI{6000}{rpm} (260/283);
0.2\,V/150\,mV at 3000: 36\,\% / 3.4\,\%. That is about the performance of the ESP32 --
consistent with 12.2, where the single biquad (BP2-win) was almost as good as the
cascade. Limits: CPU load $\approx75$--80\,\% (no room for a second filter, FFT or MEM),
no 60-2 decoder (the TDC pulse is a ground-truth helper of the simulator -- a real
implementation must decode the crank wheel, e.g.\ with the input capture unit of
Timer1, which competes with the sampling interrupt), window positions from the previous
combustion (lag at fast rpm changes), 10-bit ADC (resolution \SI{4.9}{mV}$/0.6\approx
\SI{8}{mV}$ at the sensor, sufficient for knock amplitudes $\gg$ noise), and no
floating point for diagnostics.
\end{solution}

\begin{deliverables}
12.1 on paper (block justification, numbers for b)--d)). \texttt{knockeval}: AUC and
detection tables, pooled ROC plot (BP-win, WB-win, FFT, MEM), chosen threshold with
justification. Wokwi: the confusion counts for the settings of 12.3\,d) and a
screenshot of the LEDs; a short discussion (max.\ one page) of the remaining weaknesses of
your detector.
\end{deliverables}
